% Options for packages loaded elsewhere
% Options for packages loaded elsewhere
\PassOptionsToPackage{unicode}{hyperref}
\PassOptionsToPackage{hyphens}{url}
\PassOptionsToPackage{dvipsnames,svgnames,x11names}{xcolor}
%
\documentclass[
  british,
]{article}
\usepackage{xcolor}
\usepackage{amsmath,amssymb}
\setcounter{secnumdepth}{5}
\usepackage{iftex}
\ifPDFTeX
  \usepackage[T1]{fontenc}
  \usepackage[utf8]{inputenc}
  \usepackage{textcomp} % provide euro and other symbols
\else % if luatex or xetex
  \usepackage{unicode-math} % this also loads fontspec
  \defaultfontfeatures{Scale=MatchLowercase}
  \defaultfontfeatures[\rmfamily]{Ligatures=TeX,Scale=1}
\fi
\usepackage{lmodern}
\ifPDFTeX\else
  % xetex/luatex font selection
\fi
% Use upquote if available, for straight quotes in verbatim environments
\IfFileExists{upquote.sty}{\usepackage{upquote}}{}
\IfFileExists{microtype.sty}{% use microtype if available
  \usepackage[]{microtype}
  \UseMicrotypeSet[protrusion]{basicmath} % disable protrusion for tt fonts
}{}
\makeatletter
\@ifundefined{KOMAClassName}{% if non-KOMA class
  \IfFileExists{parskip.sty}{%
    \usepackage{parskip}
  }{% else
    \setlength{\parindent}{0pt}
    \setlength{\parskip}{6pt plus 2pt minus 1pt}}
}{% if KOMA class
  \KOMAoptions{parskip=half}}
\makeatother
% Make \paragraph and \subparagraph free-standing
\makeatletter
\ifx\paragraph\undefined\else
  \let\oldparagraph\paragraph
  \renewcommand{\paragraph}{
    \@ifstar
      \xxxParagraphStar
      \xxxParagraphNoStar
  }
  \newcommand{\xxxParagraphStar}[1]{\oldparagraph*{#1}\mbox{}}
  \newcommand{\xxxParagraphNoStar}[1]{\oldparagraph{#1}\mbox{}}
\fi
\ifx\subparagraph\undefined\else
  \let\oldsubparagraph\subparagraph
  \renewcommand{\subparagraph}{
    \@ifstar
      \xxxSubParagraphStar
      \xxxSubParagraphNoStar
  }
  \newcommand{\xxxSubParagraphStar}[1]{\oldsubparagraph*{#1}\mbox{}}
  \newcommand{\xxxSubParagraphNoStar}[1]{\oldsubparagraph{#1}\mbox{}}
\fi
\makeatother


\usepackage{longtable,booktabs,array}
\newcounter{none} % for unnumbered tables
\usepackage{calc} % for calculating minipage widths
% Correct order of tables after \paragraph or \subparagraph
\usepackage{etoolbox}
\makeatletter
\patchcmd\longtable{\par}{\if@noskipsec\mbox{}\fi\par}{}{}
\makeatother
% Allow footnotes in longtable head/foot
\IfFileExists{footnotehyper.sty}{\usepackage{footnotehyper}}{\usepackage{footnote}}
\makesavenoteenv{longtable}
\usepackage{graphicx}
\makeatletter
\newsavebox\pandoc@box
\newcommand*\pandocbounded[1]{% scales image to fit in text height/width
  \sbox\pandoc@box{#1}%
  \Gscale@div\@tempa{\textheight}{\dimexpr\ht\pandoc@box+\dp\pandoc@box\relax}%
  \Gscale@div\@tempb{\linewidth}{\wd\pandoc@box}%
  \ifdim\@tempb\p@<\@tempa\p@\let\@tempa\@tempb\fi% select the smaller of both
  \ifdim\@tempa\p@<\p@\scalebox{\@tempa}{\usebox\pandoc@box}%
  \else\usebox{\pandoc@box}%
  \fi%
}
% Set default figure placement to htbp
\def\fps@figure{htbp}
\makeatother


% definitions for citeproc citations
\NewDocumentCommand\citeproctext{}{}
\NewDocumentCommand\citeproc{mm}{%
  \begingroup\def\citeproctext{#2}\cite{#1}\endgroup}
\makeatletter
 % allow citations to break across lines
 \let\@cite@ofmt\@firstofone
 % avoid brackets around text for \cite:
 \def\@biblabel#1{}
 \def\@cite#1#2{{#1\if@tempswa , #2\fi}}
\makeatother
\newlength{\cslhangindent}
\setlength{\cslhangindent}{1.5em}
\newlength{\csllabelwidth}
\setlength{\csllabelwidth}{3em}
\newenvironment{CSLReferences}[2] % #1 hanging-indent, #2 entry-spacing
 {\begin{list}{}{%
  \setlength{\itemindent}{0pt}
  \setlength{\leftmargin}{0pt}
  \setlength{\parsep}{0pt}
  % turn on hanging indent if param 1 is 1
  \ifodd #1
   \setlength{\leftmargin}{\cslhangindent}
   \setlength{\itemindent}{-1\cslhangindent}
  \fi
  % set entry spacing
  \setlength{\itemsep}{#2\baselineskip}}}
 {\end{list}}
\usepackage{calc}
\newcommand{\CSLBlock}[1]{\hfill\break\parbox[t]{\linewidth}{\strut\ignorespaces#1\strut}}
\newcommand{\CSLLeftMargin}[1]{\parbox[t]{\csllabelwidth}{\strut#1\strut}}
\newcommand{\CSLRightInline}[1]{\parbox[t]{\linewidth - \csllabelwidth}{\strut#1\strut}}
\newcommand{\CSLIndent}[1]{\hspace{\cslhangindent}#1}

\ifLuaTeX
\usepackage[bidi=basic,shorthands=off]{babel}
\else
\usepackage[bidi=default,shorthands=off]{babel}
\fi
\ifLuaTeX
  \usepackage{selnolig} % disable illegal ligatures
\fi


\setlength{\emergencystretch}{3em} % prevent overfull lines

\providecommand{\tightlist}{%
  \setlength{\itemsep}{0pt}\setlength{\parskip}{0pt}}



 


% Get Math Fonts
\usepackage{fontspec}
\usepackage{unicode-math}
\setmathfont{Latin Modern Math}
\setmathfont[range={\mathscr,\mathbfscr,\mathbb}]{XITSMath-Regular}
[    Extension = .otf,
      BoldFont = XITSMath-Bold
]
\makeatletter
\@ifpackageloaded{caption}{}{\usepackage{caption}}
\AtBeginDocument{%
\ifdefined\contentsname
  \renewcommand*\contentsname{Table of contents}
\else
  \newcommand\contentsname{Table of contents}
\fi
\ifdefined\listfigurename
  \renewcommand*\listfigurename{List of Figures}
\else
  \newcommand\listfigurename{List of Figures}
\fi
\ifdefined\listtablename
  \renewcommand*\listtablename{List of Tables}
\else
  \newcommand\listtablename{List of Tables}
\fi
\ifdefined\figurename
  \renewcommand*\figurename{Figure}
\else
  \newcommand\figurename{Figure}
\fi
\ifdefined\tablename
  \renewcommand*\tablename{Table}
\else
  \newcommand\tablename{Table}
\fi
}
\@ifpackageloaded{float}{}{\usepackage{float}}
\floatstyle{ruled}
\@ifundefined{c@chapter}{\newfloat{codelisting}{h}{lop}}{\newfloat{codelisting}{h}{lop}[chapter]}
\floatname{codelisting}{Listing}
\newcommand*\listoflistings{\listof{codelisting}{List of Listings}}
\usepackage{amsthm}
\theoremstyle{definition}
\newtheorem{exercise}{Exercise}[section]
\theoremstyle{plain}
\newtheorem{proposition}{Proposition}[section]
\theoremstyle{definition}
\newtheorem{definition}{Definition}[section]
\theoremstyle{remark}
\AtBeginDocument{\renewcommand*{\proofname}{Proof}}
\newtheorem*{remark}{Remark}
\newtheorem*{solution}{Solution}
\newtheorem{refremark}{Remark}[section]
\newtheorem{refsolution}{Solution}[section]
\makeatother
\makeatletter
\makeatother
\makeatletter
\@ifpackageloaded{caption}{}{\usepackage{caption}}
\@ifpackageloaded{subcaption}{}{\usepackage{subcaption}}
\makeatother

\newcommand{\AmenableGroup}{{\Gamma}}
\newcommand{\AmenableGroupElement}{{\gamma}}
\newcommand{\Dual}[1]{{#1}^{*}}
\newcommand{\Folner}[1][\vphantom{}]{{\Phi}_{#1}}
\newcommand{\Group}{{\Gamma}}
\newcommand{\GroupElement}{{\gamma}}
\newcommand{\HilbertSpace}{{H}}
\newcommand{\Torus}{{\mathbb{T}}}

\usepackage{bookmark}
\IfFileExists{xurl.sty}{\usepackage{xurl}}{} % add URL line breaks if available
\urlstyle{same}
% fallback for those not using the hyperref driver hyperxmp:
\makeatletter
\@ifundefined{xmpquote}{\newcommand{\xmpquote}[1]{#1}}{}
\makeatother
\hypersetup{
  pdftitle={Functional Analysis},
  pdflang={en-GB},
  colorlinks=true,
  linkcolor={blue},
  filecolor={Maroon},
  citecolor={Blue},
  urlcolor={Blue},
  pdfcreator={LaTeX via pandoc}}


\title{Functional Analysis}
\author{}
\date{2026-09-26}
\begin{document}
\maketitle


The foundational knowledge relating to Measure \& Ergodic Theory, that
has not been covered elsewhere, was provided by Jamneshan and Kreidler
(\citeproc{ref-ISem28}{2025}).

\begin{definition}[]\protect\hypertarget{def-Representation}{}\label{def-Representation}

For a Hilbert space \(\HilbertSpace\),

\begin{enumerate}
\def\labelenumi{\arabic{enumi}.}
\tightlist
\item
  A \emph{linear isometry} is a linear map
  \(U:\HilbertSpace \rightarrow \HilbertSpace\) such that
  \(||Uf||=||f||\) for all \(f\in \HilbertSpace\) (and thus injective).
\item
  A linear isometry, \(U\), is a \emph{unitary operator} if it is
  surjective (and thus bijective).
\item
  We write \(\mathscr{U}(\HilbertSpace)\) for the set of all unitary
  operators \(U:\HilbertSpace \rightarrow \HilbertSpace\).
\item
  For a group \(\Group\), we call a group homomorphism
  \(U:\Group\rightarrow\mathscr{U}(\HilbertSpace )\),
  \(\GroupElement\mapsto U_\GroupElement\) a \emph{unitary
  representation} of \(\Group\) on \(\HilbertSpace\).
\end{enumerate}

\end{definition}

\begin{definition}[]\protect\hypertarget{def-MarkovSubLattive}{}\label{def-MarkovSubLattive}

Let \(J:(Y,S)\rightarrow(X,T)\) be an extension of measure-preserving
systems. Then \(E=U_J(\text{L}^2(Y))\) is an \emph{invariant Markov
sub-lattice} of \(\text{L}^2(X)\), i.e.,

\begin{enumerate}
\def\labelenumi{\arabic{enumi}.}
\tightlist
\item
  \(E\) is a closed linear subspace of \(\text{L}^2(X)\),
\item
  \(\boldsymbol{1}\in E\),
\item
  \(|f|,\ \text{Re}(f),\ \text{Im}(f)\in E\) for every \(f\in E\), and
\item
  \(U_{T_\GroupElement}f\in E\) for every \(f\in E\) and
  \(\GroupElement\in\Group\).
\end{enumerate}

\end{definition}

\begin{definition}[]\protect\hypertarget{def-ContractionSemigroup}{}\label{def-ContractionSemigroup}

Let \(\HilbertSpace\) be a Hilbert space. \(\mathscr{L}(\HilbertSpace)\)
is the space of all bounded linear operators from \(\HilbertSpace\) to
\(\HilbertSpace\). \footnote{Note,
  \(\mathscr{U}(\HilbertSpace)\subseteq\mathscr{L}(\HilbertSpace)\) and
  \(||U||=1\) when \(U\in\mathscr{U}(\HilbertSpace)\), so
  \(\mathscr{U}(\HilbertSpace)\) is a contraction semigroup.}

A family \(\mathscr{S}\subseteq\mathscr{L}(\HilbertSpace)\) is a
\emph{semigroup (of operators)} if \(UV\in\mathscr{S}\) for all
\(U,V\in\mathscr{S}\). It is a \emph{contraction semigroup} if, in
addition, \(||U||\leq1\) for all \(U\in\mathscr{S}\).

We call
\[\text{fix}(\mathscr{S}):=\bigcap_{U\in\mathscr{S}}\text{fix}(U)=\{f\in\HilbertSpace\mid Uf=f\text{ for every }U\in\mathscr{S}\}\]
the \emph{fixed space} of \(\mathscr{S}\).\footnote{\(\text{fix}(\mathscr{S})\subseteq \HilbertSpace_\text{ds}\).}\footnote{Meeting
  Notes: 0-dimensional. Next most simple case would be scaling
  (1-dimensional), and then rotations on a place (2-dimensional).}

The \emph{closed convex hull} \(\overline{\text{co}}\ A\) is the closure
of the set of all convex combinations of elements of \(A\).

The \emph{closed linear hull} \(\overline{\text{lin}}\ A\) is the
closure of the set of all linear combinations of elements of \(A\).

\end{definition}

\begin{definition}[]\protect\hypertarget{def-Torus}{}\label{def-Torus}

We write \(\Torus:=\{z\in\mathbb{C}\mid |z|=1\}\) for the multiplicative
group of complex numbers of modulus one, and call this the \emph{torus}.

\end{definition}

\begin{definition}[]\protect\hypertarget{def-DualGroup}{}\label{def-DualGroup}

Let \(U:\AmenableGroup\rightarrow\mathscr{U}(\HilbertSpace)\) be a
unitary representation of a discrete abelian amenable group
\(\AmenableGroup\).

A group homomorphism \(\chi:\AmenableGroup\rightarrow\Torus\), where
\(\mathbb{T}=\{z\in\mathbb{C}\mid|z|=1\}\), is called a
\emph{character}. The \emph{dual group} \(\Dual{\AmenableGroup}\) of
\(\AmenableGroup\) is the set of all such characters equipped with the
multiplication given by
\((\chi_1\chi_2)(\AmenableGroupElement):=\chi_1(\AmenableGroupElement)\chi_2(\AmenableGroupElement)\)
for \(\AmenableGroupElement\in\AmenableGroup\) and
\(\chi_1,\chi_2\in\Dual{\AmenableGroup}\).\footnote{Commutator subgroup
  and abelianisation?}

\end{definition}

\begin{definition}[]\protect\hypertarget{def-DiscreteSpectrum}{}\label{def-DiscreteSpectrum}

Let
\[M_f:=\text{lin}\{U_\AmenableGroupElement f\mid\AmenableGroupElement\in\AmenableGroup \}. \]
We call a subset \(M\subseteq \HilbertSpace\) an \emph{invariant
finite-dimensional subspace} if \(M_f\) is finite-dimensional and
\(U_\AmenableGroupElement f\in M\) for all \(f\in M\) and
\(\AmenableGroupElement\in\AmenableGroup\).

The closure
\[\HilbertSpace_\text{ds}=\overline{\bigcup\{M\subseteq \HilbertSpace\mid M\text{ invariant finite-dimensional subspace} \}}\subseteq \HilbertSpace \]
is called the \emph{discrete spectrum part} of \(U\).\footnote{Meeting
  Notes: In the \(\AmenableGroup=\mathbb{Z}\) case, then
  \(\HilbertSpace_\text{ds}\) can be broken up into one-dimensional
  objects that corresponds to an eigenfunction.}

\end{definition}

\begin{definition}[{Jamneshan and Kreidler, 2025, Definition
5.1.8}]\protect\hypertarget{def-EigenSpace}{}\label{def-EigenSpace}

Let \(U:\Group\rightarrow\mathscr{U}(\HilbertSpace)\) be a unitary
representation of an abelian group \(\Group\). For a character
\(\chi\in\Dual{\Group}\), we call \begin{align*}
  \text{ker}(\chi-U):&=\{f\in\HilbertSpace\mid U_\GroupElement f=\chi(\GroupElement)f\text{ for every }\GroupElement\in\Group\}
  \\&\bigcap_{\GroupElement\in\Group}\text{ker}(\chi(\GroupElement)\text{Id}_\HilbertSpace-U_\GroupElement)
\end{align*} the \emph{eigenspace} of \(U\) with respect to \(\chi\),
and elements \(f\in\text{ker}(\chi-U)\setminus\{0\}\) are called
\emph{eigenvectors} with respect to \(\chi\). The set
\(\sigma_p(U):=\{\chi\in\Dual{\Group}\mid\text{ker}(\chi-U)\neq\{0\}\}\)
is the \emph{point spectrum} of \(U\) and its elements are called the
\emph{eigenvalues} of \(U\).

\end{definition}

\begin{proposition}[{cf. Jamneshan and Kreidler, 2025, Proposition
5.1.10}]\protect\hypertarget{prp-POrthogonalEigenspaces}{}\label{prp-POrthogonalEigenspaces}

The eigenspaces \(\text{ker}(\chi-U)\) for \(\chi\in\Dual{\Group}\) are
pairwise orthogonal. For \(M\subseteq \HilbertSpace\), then the
following are equivalent:

\begin{enumerate}
\def\labelenumi{\arabic{enumi}.}
\tightlist
\item
  \(M\) is an irreducible invariant finite-dimensional subspace.
\item
  \(M\) is an invariant linear subspace which is at most
  one-dimensional.
\item
  There is \(\chi\in\Dual{\Group}\) and \(f\in\text{ker}(\chi-U)\), such
  that \(M=\mathbb{C}\cdot f\).
\end{enumerate}

\end{proposition}

\begin{definition}[]\protect\hypertarget{def-StronglyContinuous}{}\label{def-StronglyContinuous}

A unitary representation
\(U:\AmenableGroup\rightarrow\mathscr{U}(\HilbertSpace)\) of a
topological group \(\AmenableGroup\) is \emph{strongly continuous} if
\(\AmenableGroup\rightarrow \HilbertSpace\), \(x\mapsto U_xf\) is
continuous for every \(f\in \HilbertSpace\).

\end{definition}

\begin{definition}[]\protect\hypertarget{def-RegularRepresentation}{}\label{def-RegularRepresentation}

For every discrete group \(\AmenableGroup\), the maps \begin{align}
    L:\AmenableGroup\rightarrow\mathscr{U}(\text{L}^2(\AmenableGroup)),&\qquad x\mapsto L_x
    \\ R:\AmenableGroup\rightarrow\mathscr{U}(\text{L}^2(\AmenableGroup)),&\qquad x\mapsto R_x
\end{align} with \(L_xf:=f\circ l_{x^{-1}}\) and
\(R_xf:=f\circ r_{x^{-1}}\) for \(f\in\text{L}^2(\AmenableGroup)\) and
\(x\in\AmenableGroup\) are strongly continuous unitary representations.
We call \(L\) and \(R\) the \emph{left} and \emph{right regular
representation} of \(\AmenableGroup\), respectively.\footnote{As
  \(\AmenableGroup\) is discrete, then any map
  \(\AmenableGroup\rightarrow \HilbertSpace\) is continuous for any
  topological space \(\HilbertSpace\). Thus,
  \(\AmenableGroup\rightarrow \text{L}^2(\AmenableGroup)\) defined by
  \(\AmenableGroupElement\mapsto L_\AmenableGroupElement f\) or
  \(\AmenableGroupElement\mapsto R_\AmenableGroupElement f\) is always
  continuous. Hence, \(L\) and \(R\) are strongly continuous unitary
  representations.}

\end{definition}

Assume that \(U:\AmenableGroup\rightarrow\mathscr{U}(\HilbertSpace)\) is
a strongly continuous unitary representation of a discrete group
\(\AmenableGroup\). Let further \(P\) be the orthogonal projection onto
the fixed space \(\text{fix}(U(\AmenableGroup))\). Then
\[(Pf\mid g)=\lim_{N\rightarrow\infty}\frac{1}{|\Folner[N]|}\sum_{\AmenableGroupElement\in\Folner[N]}(U_\AmenableGroupElement f\mid g)\]
for all \(f,g\in\HilbertSpace\).

If \(U:\AmenableGroup\rightarrow\mathscr{U}(\HilbertSpace)\) is a
strongly continuous unitary representation of a compact group
\(\AmenableGroup\), then \(U\) has discrete spectrum.\footnote{This
  \(\AmenableGroup\) would be the compact group behind the factor
  subsystem of the larger system. For example,
  \(\AmenableGroup=\mathbb{Z}\) for the irrational \(\alpha\) group
  rotation would give us compact \(G=\{z\in\mathbb{C}\mid|z|=1\}\). The
  homomorphism from \(\AmenableGroup\rightarrow \Group\) is
  \(n\mapsto e^{2\pi i\alpha n}\).}

\begin{exercise}[]\protect\hypertarget{exr-AmenableAbstractMeanErgodic}{}\label{exr-AmenableAbstractMeanErgodic}

Show that
\[\HilbertSpace=\HilbertSpace_\text{ds}\oplus \HilbertSpace_\text{wm} \]
and what happens to \(\HilbertSpace_\text{ds}\) and
\(\HilbertSpace_\text{wm}\) when it is averaged by the Abstract Mean
Ergodic Theorem for amenable groups.\footnote{(Proposition
  5.3.3)\[\HilbertSpace_\text{wm}=\left\{f\in \HilbertSpace\mid \lim_{N\rightarrow\infty}\frac{1}{|\Folner[N]|}\sum_{\varphi\in\Folner[N]}|(U_\varphi f\mid f)|^2=0  \right\}. \]}

\end{exercise}

\begin{proposition}[]\protect\hypertarget{prp-DSInvariantMarkov}{}\label{prp-DSInvariantMarkov}

Let \((X,T)\) be a measure-preserving system. Then
\(\text{L}^2(X)_\text{ds}\) is an invariant Markov sublattice of
\(\text{L}^2(X)\).

\end{proposition}

\begin{exercise}[{Jamneshan and Kreidler, 2025, Lemma
6.1.12}]\protect\hypertarget{exr-ISem28exr6112}{}\label{exr-ISem28exr6112}

Assume that \(U:\AmenableGroup\rightarrow\mathscr{U}(\HilbertSpace)\) is
a strongly continuous unitary representation of a discrete group
\(\AmenableGroup\). Let further \(P\) be the orthogonal projection onto
the fixed space \(\text{fix}(U(\AmenableGroup))\). Then
\((Pf\mid g)=\lim_{N\rightarrow\infty}\frac{1}{|\Folner[N]|}\sum_{\AmenableGroupElement\in\Folner[N]}(U_\AmenableGroupElement f\mid g)\)
for all \(f,g\in\HilbertSpace\).

\end{exercise}

\protect\phantomsection\label{refs}
\begin{CSLReferences}{1}{0}
\bibitem[\citeproctext]{ref-ISem28}
Jamneshan, A. and Kreidler, H. (2025). \emph{ISem 28: Ergodic structure
theory and applications}.

\end{CSLReferences}




\end{document}
